\section*{Chapter \ref{ch:s2:manipulative-strategies-and-how-they-are-caught} --- Manipulative Strategies and How They Are Caught}

\begin{solution}{exo:s2:manipulative-strategies-and-how-they-are-caught:1}
$24\,814 \times 0.005 \approx 124$ orders traded; the other 24\,690 were cancelled.
\end{solution}

\begin{solution}{exo:s2:manipulative-strategies-and-how-they-are-caught:2}
$0.01 \times 11\,500 = 115$ legitimate account-days flagged, against 77 true episodes: fewer than half the alerts are true.
\end{solution}

\begin{solution}{exo:s2:manipulative-strategies-and-how-they-are-caught:3}
About \$38.6 million.
\end{solution}

\begin{solution}{exo:s2:manipulative-strategies-and-how-they-are-caught:4}
Cancelling a lot is what market makers do: they quote continuously and replace their quotes as prices move, and cancel more than the spoofers in the
sample. What distinguishes a spoofer is cancelling one kind of order while trading another, which the ratio averages away.
\end{solution}

\begin{solution}{exo:s2:manipulative-strategies-and-how-they-are-caught:5}
Index funds and other investors must trade at the close, because their benchmarks are closing prices; heavy closing volume is their normal
behaviour. The Optiver record adds what matters: a position priced at the settlement opposite to the closing trades, and a move that reverses.
\end{solution}

\begin{solution}{exo:s2:manipulative-strategies-and-how-they-are-caught:6}
Almost every order in modern markets is cancelled, and legitimate traders cancel when prices move or after fills; a rule based on cancellation alone
would prohibit market making. The statute and the Market Abuse Regulation target orders placed to be cancelled, or to mislead, which courts infer
from design and patterns, as in Coscia.
\end{solution}

\begin{solution}{exo:s2:manipulative-strategies-and-how-they-are-caught:7}
63\% at a 0.1\% false-positive rate (about 11 legitimate alerts) and 100\% at 10\% (about 1\,150). Catching the last third of episodes costs a hundred
times more alerts: the threshold is set by how many alerts investigators can read.
\end{solution}

\begin{solution}{exo:s2:manipulative-strategies-and-how-they-are-caught:8}
The order-to-trade ratio caught none of the planted spoofers and flagged market makers instead. Having no alerts from a detector that cannot see the
pattern says nothing about whether the pattern is there.
\end{solution}

\begin{solution}{pb:s2:manipulative-strategies-and-how-they-are-caught:1}
\begin{enumerate}
\item \omterm{def:s2:manipulative-strategies-and-how-they-are-caught:spoofing}{Spoofing} is bidding or offering with the intent to cancel before execution; \omterm{def:s2:manipulative-strategies-and-how-they-are-caught:layering}{layering} is \omterm{def:s2:manipulative-strategies-and-how-they-are-caught:spoofing}{spoofing} with several orders at successive levels kept
away from the best price.
\item Small and large orders on opposite sides; large orders cancelled after time, a small fill or a single large fill; 24\,814 large orders of which
0.5\% traded, small orders about 52\% filled; over 450\,000 large orders and \$1.4 million.
\item A modified platform \omterm{def:s2:manipulative-strategies-and-how-they-are-caught:layering}{layering} four to six large sell orders three or four levels from the best ask, on over 400 days, profits over \$40 million,
and close to \$200 million of pressure before the Flash Crash; \$38.6 million by consent in 2016.
\item \omterm{def:s2:manipulative-strategies-and-how-they-are-caught:spoofing}{Spoofing} as defined in 7~U.S.C.~6c(a)(5)(C); the regulation lists false signals, trading at the close, orders that initiate trends, \omterm{def:s2:manipulative-strategies-and-how-they-are-caught:benchmark}{benchmark
manipulation}, and indicators such as orders removed before execution.
\item \omterm{def:s2:manipulative-strategies-and-how-they-are-caught:ignition}{Momentum ignition}: orders intended to start or exaggerate a trend; \omterm{def:s2:manipulative-strategies-and-how-they-are-caught:close}{marking the close}: heavy trading at the close to move a closing price.
\item Traders offset large Trading at Settlement positions by trading against them in the closing minutes; 19 attempts, 5 successful; \$14 million.
\item False inputs or coordinated trading to move a benchmark; four banks pleaded guilty in 2015 to manipulating fixes, with fines over \$2.5 billion.
\item Three purported market makers and nine individuals wash-traded crypto assets for promoters, with bots generating billions of dollars of volume a
day.
\item Market makers, quote refreshers, deep-book providers and directional traders, each tripping one detector; planted spoofers following the
Coscia pattern.
\item \omterm{def:s2:manipulative-strategies-and-how-they-are-caught:spoofing}{Spoofing}: 0\%, 37\%, 51\% and 77\% caught; close: 1\%, 21\%, 61\%; wash: 74\% and 97\%, each at 1\% false positives.
\item Each feature has innocent causes that rarely occur together; the spoofer's pattern needs both.
\item Unsettled: no court ruling was found; it is avoided by private transaction channels.
\item At a 1\% false-positive rate the combined \omterm{def:s2:manipulative-strategies-and-how-they-are-caught:spoofing}{spoofing} detector catches 77\% and the order-to-trade ratio 0\%, flagging market makers; the close
detector catches 61\% and flags 3.3\% of index funds; the self-match share catches 97\% and flags 5.0\% of multi-algorithm firms.
\item It is ranked and investigated: the orders behind it are read, and intent is looked for in design, communications and records.
\item On legitimate look-alikes, at the operating threshold, counting alerts per investigator, with episodes from real cases where possible.
\item Mix order sizes and timings; the more the large orders look like genuine orders, the more often they trade, which is exactly the risk the
manipulation was designed to avoid.
\item They describe the observable patterns that define the offence and that courts accepted as evidence of intent.
\item The cancellation-pattern surveillance detector.
\item Mark-outs price a client's information; surveillance asks whether an account's orders were meant to trade at all.
\item The market maker's orders are meant to trade; the spoofer's are meant not to.
\end{enumerate}
\end{solution}

\begin{solution}{iq:s2:manipulative-strategies-and-how-they-are-caught:1}
By showing that orders are genuine: they trade at similar rates across sizes and sides given their distance from the touch, cancellations follow price
moves and risk limits rather than fills on the other side, and the code's cancellation rules have legitimate purposes, documented before trading.
\end{solution}

\begin{solution}{iq:s2:manipulative-strategies-and-how-they-are-caught:2}
Per account and day, compare fill rates of large and small orders and the timing of large cancellations relative to fills on the other side; score
their combination; test on labelled cases and on legitimate strategies chosen to look similar; report the catch rate at the operating false-positive
rate.
\end{solution}

\begin{solution}{iq:s2:manipulative-strategies-and-how-they-are-caught:3}
Trade for genuine reasons (hedging, client orders) in ways that do not aim to move the settlement; avoid concentrating trading against the position in
the closing window; document the reasons; check the firm's rules and the venue's.
\end{solution}

\begin{solution}{iq:s2:manipulative-strategies-and-how-they-are-caught:4}
Pre-trade review of algorithms' cancellation and order-size logic; self-match prevention; surveillance of the firm's own order logs with the detectors
above; limits on closing-window trading against settlement-priced positions; communications surveillance; training and escalation.
\end{solution}

\begin{solution}{iq:s2:manipulative-strategies-and-how-they-are-caught:5}
Every order event (new, modify, cancel, fill) with exchange and local timestamps, account and beneficial owner, algorithm and version, price, size,
side and venue, and the reason for cancellations where the system knows it; kept complete and unaltered.
\end{solution}

\begin{solution}{iq:s2:manipulative-strategies-and-how-they-are-caught:6}
$77/(77 + 115) = 40.1\%$ of alerts are true. With ten episodes instead of 100 and the same threshold, $7.7/(7.7 + 115) = 6.3\%$: rarer offences make
every detector's alerts mostly false.
\end{solution}
